\documentclass[11pt]{article}

\usepackage[margin=1in]{geometry}
\usepackage{amsmath,amssymb,amsthm}
\usepackage{graphicx}
\usepackage[hidelinks]{hyperref}

\newtheorem{theorem}{Theorem}
\newtheorem{corollary}{Corollary}

\newcommand{\D}{\mathbb D}
\newcommand{\N}{\mathbb N}

\title{Bohr-lift zero incomparability}
\author{Run \texttt{fa\_banach\_001}}
\date{2026-06-28}

\begin{document}
\maketitle

\section*{Source and Target}

The source paper is Jeet Sampat's arXiv:2409.10224 paper
\cite{Sampat2024}, \emph{Shift-cyclicity in analytic function spaces}.  The
target is Problem 2.23 on PDF page 20.  The paper defines the Bohr transform
\(\mathfrak B:H^2_0(\D)\to H^2(\D^\infty_2)\) by sending the coefficient of
\(z^n\) to the monomial indexed by the prime factorization of \(n\).  The
problem asks how \(f(\D)\) and \(\mathfrak Bf(\D^\infty_2)\) are related and,
in particular, asks for a classification of those \(f\in H^2_0(\D)\) for which
\(\mathfrak Bf\) is non-vanishing on \(\D^\infty_2\).

\begin{center}
\includegraphics[width=\textwidth,height=0.78\textheight,keepaspectratio]{figures/open_problem_crop.png}
\end{center}

\section*{Idea of the Proof}

The first three positive indices already split under the Bohr transform:
\(1\) has the empty prime factorization, while \(2\) and \(3\) become the first
two coordinate variables.  Thus the ordinary one-variable cubic
\[
  f(z)=az+bz^2+cz^3
\]
is turned into the affine two-variable function
\[
  \mathfrak Bf(w)=a+bw_1+cw_2.
\]
Nonvanishing on the Hilbert multidisk is therefore controlled by whether the
disk centered at \(a\) with radius \(|b|+|c|\) reaches the origin.  This gives
a complete degree-three criterion and a simple example whose ordinary disk
zeros do not predict Bohr-lift zeros.  A second finite example, using the
factorization \((1-w_1)(1-w_2)\), shows the reverse failure: Bohr-lift
nonvanishing does not imply ordinary disk nonvanishing away from the forced
zero at the origin.

\begin{theorem}
Let
\[
  f(z)=az+bz^2+cz^3\in H^2_0(\D).
\]
Then \(\mathfrak Bf\) is non-vanishing on \(\D^\infty_2\) if and only if
\[
  a\neq 0\quad\text{and}\quad |a|\ge |b|+|c|.
\]
\end{theorem}

\begin{proof}
Because \(1\) has the empty prime factorization, while \(2=p_1\) and
\(3=p_2\), the Bohr transform is
\[
  \mathfrak Bf(w)=a+bw_1+cw_2,\qquad w=(w_1,w_2,\ldots)\in\D^\infty_2.
\]
If \(a=0\), then \(\mathfrak Bf(0)=0\), so nonvanishing is impossible.

Assume \(a\neq 0\) and \(|a|\ge |b|+|c|\).  For every
\(w\in \D^\infty_2\),
\[
  |bw_1+cw_2| < |b|+|c|\le |a|,
\]
where the first inequality is strict unless \(b=c=0\), in which case the
claim is immediate.  Hence \(a+bw_1+cw_2\neq 0\).

Conversely, suppose \(|a|<|b|+|c|\).  If \(a=0\), the zero at \(w=0\) has
already been noted.  Otherwise choose \(0\le r_1,r_2<1\) such that
\[
  |a|=|b|r_1+|c|r_2,
\]
with the evident convention that a term with zero coefficient is omitted.
Choose phases for \(w_1,w_2\) so that \(bw_1\) and \(cw_2\) both point in the
direction \(-a\), and set all remaining coordinates to \(0\).  Then
\(w\in\D^\infty_2\) and \(a+bw_1+cw_2=0\).  This proves the criterion.
\end{proof}

\begin{corollary}
The function \(f_0(z)=z+z^2+z^3\) has no zeros in \(\D\setminus\{0\}\), but
\(\mathfrak Bf_0\) vanishes on \(\D^\infty_2\).
\end{corollary}

\begin{proof}
We have \(f_0(z)=z(1+z+z^2)\), and the zeros of \(1+z+z^2\) are the primitive
third roots of unity, which lie on \(\partial\D\).  On the other hand,
\[
  \mathfrak Bf_0(w)=1+w_1+w_2,
\]
so \(\mathfrak Bf_0(-1/2,-1/2,0,\ldots)=0\).  The point
\((-1/2,-1/2,0,\ldots)\) belongs to \(\D^\infty_2\).
\end{proof}

\begin{theorem}[Zero incomparability]
The two zero conditions
\[
  0\notin f(\D\setminus\{0\})
  \qquad\text{and}\qquad
  0\notin \mathfrak Bf(\D^\infty_2)
\]
are logically independent, even for polynomials in \(H^2_0(\D)\).
\end{theorem}

\begin{proof}
The preceding corollary gives one direction: \(f_0(z)=z+z^2+z^3\) has no zero
in \(\D\setminus\{0\}\), but its Bohr lift vanishes at
\((-1/2,-1/2,0,\ldots)\).

For the reverse direction, set
\[
  g(z)=z-z^2-z^3+z^6.
\]
The prime factorizations \(1=1\), \(2=p_1\), \(3=p_2\), and \(6=p_1p_2\)
give
\[
  \mathfrak Bg(w)=1-w_1-w_2+w_1w_2=(1-w_1)(1-w_2).
\]
Since \(|w_1|<1\) and \(|w_2|<1\) for \(w\in\D^\infty_2\), neither factor
can vanish.  Hence \(\mathfrak Bg\) is non-vanishing on \(\D^\infty_2\).

On the other hand \(g(z)=z h(z)\), where
\[
  h(z)=1-z-z^2+z^5.
\]
Now \(h(0)=1\), while
\[
  h(3/4)=1-\frac34-\frac{9}{16}+\frac{243}{1024}
        =-\frac{77}{1024}<0.
\]
By the intermediate value theorem, \(h\) has a real zero in \((0,3/4)\).  Thus
\(g\) has a zero in \(\D\setminus\{0\}\), although \(\mathfrak Bg\) is
non-vanishing on \(\D^\infty_2\).
\end{proof}

\section*{Verification Notes}

These are finite-dimensional calculations inside the Bohr transform.  No
computer search was needed.  The first proof uses only the prime
factorizations \(1=1\), \(2=p_1\), and \(3=p_2\), and the fact that the first
two coordinates of \(\D^\infty_2\) can be chosen independently in the unit
disk.  The reverse example additionally uses \(6=p_1p_2\).

The result is partial.  It does not classify all \(f\in H^2_0(\D)\).  It shows
that even for finite positive-index polynomials, Bohr-lift nonvanishing is a
several-variable stability condition rather than ordinary nonvanishing of the
one-variable function away from the forced zero at the origin.  In particular,
the two zero conditions are independent.

\section*{Novelty and Literature Check}

The lightweight run indexes were searched for arXiv:2409.10224 and core
phrases including \texttt{Bohr transform}, \texttt{Hilbert multidisk}, and
\texttt{non-vanishing}; no duplicate packet or attempt was found.  Web/arXiv
searches on 2026-06-28 for matching Bohr-transform nonvanishing phrases,
including \(H^2_0\), found no exact duplicate of these finite-polynomial
obstructions.  The result should therefore be treated as a small new partial
observation, subject to human review.

\section*{Human Review Recommendation}

Check the source's indexing convention for \(H^2_0(\D)\) and the Bohr
transform.  In particular, the coefficient of \(z^1\) becomes the constant
term of \(\mathfrak Bf\).  Once this convention is confirmed, the results are
elementary and likely valid obstructions for Problem 2.23.

\begin{thebibliography}{9}

\bibitem{Sampat2024}
Jeet Sampat.
\newblock \emph{Shift-cyclicity in analytic function spaces}.
\newblock arXiv:2409.10224, 2024.

\end{thebibliography}

\end{document}
